\documentclass[level=2]{kidscode}
\newcommand{\letter}[1]{\par\needspace{4\baselineskip}\medskip\noindent{\Large\bfseries\color{kcaccent}#1}\par\nopagebreak}
\newenvironment{terms}{\begin{description}[leftmargin=0pt,itemsep=3pt,font=\bfseries]}{\end{description}}
\newcommand{\term}[3]{\item[#1] #2\ifx&#3&\else\ \textcolor{kcgray}{(week #3)}\fi}

\begin{document}
\kctitle{Glossary}{Level 2 words from A to Z, in plain language}

The week in brackets tells you where the word is taught, so you can look back in your workbook.
Words in \texttt{this font} are things you type in Python.

\begin{multicols}{2}
\raggedright

\letter{A}
\begin{terms}
\term{\texttt{and}}{Joins two conditions. Both must be True: \texttt{score > 5 and lives > 0}.}{1}
\term{\texttt{append}}{Adds an item to the end of a list: \texttt{names.append("Zoe")}.}{3}
\term{argument}{The value you put in the brackets when you call a function. In \texttt{area(3, 4)} the arguments are 3 and 4.}{2}
\term{\texttt{assert}}{Checks that something is True. If it is not, the test fails and shows your message.}{8}
\term{attribution}{Saying who made the art, sound or code you used. Giving credit.}{9}
\end{terms}

\letter{B}
\begin{terms}
\term{bool}{The type for True or False.}{1}
\term{\texttt{break}}{Leaves a loop straight away.}{4}
\term{bug}{A mistake in a program. Every programmer makes them; finding them is part of the fun.}{}
\end{terms}

\letter{C}
\begin{terms}
\term{Caesar cipher}{A secret code that shifts every letter along the alphabet by the same amount.}{5}
\term{code review}{Reading someone's code to ask questions, find bugs and say what is good.}{10}
\term{collision}{When two things in a game touch. We check it with \texttt{distance}.}{9}
\term{comment}{A note for people, after a \texttt{\#}. Python ignores it.}{2}
\term{compound condition}{Two or more conditions joined with \texttt{and}, \texttt{or} or \texttt{not}.}{1}
\end{terms}

\letter{D}
\begin{terms}
\term{data type}{The kind of value: int, float, str or bool. Also just called a type.}{1}
\term{debugging}{Finding and fixing bugs with a strategy: read the traceback, print-debugging, rubber duck.}{6}
\term{decomposition}{Breaking a big problem into small parts you can solve one at a time.}{10}
\term{dictionary (dict)}{A collection that stores values under names (keys): \texttt{\{"lives": 3\}}.}{5}
\term{docstring}{A comment in triple quotes as the first line of a function, saying what it does.}{2}
\term{driver}{In pair programming, the person at the keyboard.}{}
\end{terms}

\letter{E}
\begin{terms}
\term{edge case}{A test at the very edge of what is allowed: zero, an empty list, a wall.}{8}
\term{encryption}{Scrambling a message so only people with the key can read it.}{5}
\term{event}{Something that happens, like a key press, a click or a timer going off.}{9}
\term{event handler}{The function that runs when an event happens.}{9}
\end{terms}

\letter{F}
\begin{terms}
\term{f-string}{A string with an \texttt{f} in front that fills in values: \texttt{f"Score: \{score\}"}.}{1}
\term{file}{Data saved on the computer, so it is still there after the program stops.}{7}
\term{\texttt{FileNotFoundError}}{The error you get when you try to read a file that does not exist (yet).}{7}
\term{float}{The type for decimal numbers like 3.5.}{1}
\term{flowchart}{A drawing of a program's steps with boxes and arrows.}{10}
\term{function}{A named block of code you can use again. Made with \texttt{def}.}{2}
\end{terms}

\letter{G}
\begin{terms}
\term{game loop}{The part of a game that runs again and again: move, check, draw.}{4, 9}
\term{game state}{All the values that describe the game right now: score, lives, positions.}{9}
\end{terms}

\letter{I}
\begin{terms}
\term{index}{The position of an item in a list. The first item has index 0.}{3}
\term{\texttt{IndexError}}{The error you get when you ask for an index that is not in the list.}{3}
\term{infinite loop}{A loop that never stops. Press the red Stop button.}{4}
\term{input validation}{Checking that what the user typed makes sense, and asking again if not.}{4}
\term{int}{The type for whole numbers like 7 or $-3$.}{1}
\term{item}{One value inside a list. Also called an element.}{3}
\end{terms}

\letter{K}
\begin{terms}
\term{key}{The name you use to look up a value in a dictionary.}{5}
\term{\texttt{KeyError}}{The error you get when a key is not in the dictionary. Check the spelling!}{5}
\end{terms}

\letter{L}
\begin{terms}
\term{\texttt{len}}{Tells you how many items are in a list (or letters in a string).}{3}
\term{list}{A collection of values in order: \texttt{["red", "blue"]}.}{3}
\term{\texttt{listen}}{Makes the turtle window pay attention to the keyboard.}{9}
\end{terms}

\letter{M}
\begin{terms}
\term{mode}{How a file is opened: \texttt{"r"} read, \texttt{"w"} write (starts empty), \texttt{"a"} append (adds to the end).}{7}
\end{terms}

\letter{N}
\begin{terms}
\term{navigator}{In pair programming, the person who reads the task, thinks ahead and spots bugs.}{}
\term{nested loop}{A loop inside another loop. Great for grids and pixel art.}{4}
\term{normal case}{A test with ordinary, everyday values.}{8}
\term{\texttt{not}}{Turns True into False and False into True.}{1}
\end{terms}

\letter{O}
\begin{terms}
\term{\texttt{onkey}}{Tells turtle which function to run when a key is pressed.}{9}
\term{\texttt{ontimer}}{Tells turtle to run a function after some milliseconds.}{9}
\term{open license}{Permission from a creator that lets others use their work, often if they give credit.}{9}
\term{\texttt{or}}{Joins two conditions. At least one must be True.}{1}
\end{terms}

\letter{P}
\begin{terms}
\term{parameter}{A name in a function's brackets that receives a value: \texttt{def area(width, height)}.}{2}
\term{personal information}{Anything that tells who or where you are: real name, address, school, birthday, photos.}{1}
\term{playtest}{Letting someone try your game while you watch and take notes.}{11}
\term{print-debugging}{Adding \texttt{print} lines to see what your variables hold.}{6}
\term{pseudocode}{Planning a program in plain words before writing Python.}{10}
\term{pytest}{A tool that finds your \texttt{test\_} functions and runs them all.}{8}
\end{terms}

\letter{R}
\begin{terms}
\term{\texttt{random.choice}}{Picks one random item from a list.}{3}
\term{\texttt{return}}{Hands a value back from a function to the code that called it.}{2}
\term{return value}{The value a function hands back.}{2}
\term{rubber duck debugging}{Explaining your code line by line out loud (to a toy or a friend) until you spot the bug.}{6}
\end{terms}

\letter{S}
\begin{terms}
\term{scope}{Where a variable can be used. Variables made inside a function stay inside it.}{2}
\term{\texttt{split}}{Cuts a string into a list of words: \texttt{"a b".split()} gives \texttt{["a", "b"]}.}{7}
\term{str (string)}{The type for text, written in quotes: \texttt{"hello"}.}{1}
\term{\texttt{strip}}{Removes spaces and new-line characters from both ends of a string.}{7}
\end{terms}

\letter{T}
\begin{terms}
\term{test}{A small function that checks your code works, using \texttt{assert}.}{8}
\term{trace table}{A table on paper where you write how variables change, step by step.}{6}
\term{traceback}{The error report Python shows. Read the last line first.}{6}
\term{\texttt{TypeError}}{The error you get when you mix types that do not fit, like \texttt{"5" + 5}.}{1}
\end{terms}

\letter{V}
\begin{terms}
\term{value}{A piece of data, like 7, \texttt{"hi"} or True. In a dictionary, what is stored under a key.}{1, 5}
\term{variable}{A name that stores a value. Choose clear names like \texttt{lives}, not \texttt{x}.}{1}
\end{terms}

\letter{W}
\begin{terms}
\term{\texttt{while} loop}{Repeats while a condition is True. Something inside must change so it can stop.}{4}
\term{\texttt{with}}{Opens a file and closes it again safely when the block ends.}{7}
\end{terms}

\end{multicols}

\subsection*{More words from your workbook}
\begin{multicols}{2}
\begin{description}[leftmargin=0pt,itemsep=3pt,font=\bfseries]
\item[.items()] Gives every key together with its value, so a loop can use both. {\small\color{kcgray}(week 5)}
\item[accessibility] designing so that everyone can use it, including people who see colors differently or read more slowly. {\small\color{kcgray}(week 11)}
\item[accessible] Designed so that everyone can use it, including people who see colors differently. {\small\color{kcgray}(week 3)}
\item[bug list] a list of problems found in testing, sorted so the most important ones get fixed first. {\small\color{kcgray}(week 11)}
\item[bug report] A short note: what went wrong, where, and what you changed. {\small\color{kcgray}(week 6)}
\item[core loop] the part of a game or app that repeats over and over; the smallest version that works. {\small\color{kcgray}(week 10)}
\item[data privacy] Deciding what information about people is saved, who can see it, and keeping it safe. {\small\color{kcgray}(week 7)}
\item[design decision] a choice you made on purpose, and the reason for it. {\small\color{kcgray}(week 12)}
\item[encode / decode] Turn a message into secret code / turn it back again. {\small\color{kcgray}(week 5)}
\item[expected vs actual] What the test says should happen, and what really happened. {\small\color{kcgray}(week 8)}
\item[frame] One tiny step of a game: everything moves a little, then it is drawn. {\small\color{kcgray}(week 9)}
\item[gallery walk] visiting projects one by one, like in a museum, and leaving feedback. {\small\color{kcgray}(week 12)}
\item[license / attribution] The rules for using someone's work / giving them credit for it. {\small\color{kcgray}(week 9)}
\item[milestone] a checkpoint on the way to the finished project, like M2 \enquote{core loop works}. {\small\color{kcgray}(week 10)}
\item[None] Python's word for \enquote{nothing}. A function with no \texttt{return} gives back \texttt{None}. {\small\color{kcgray}(week 2)}
\item[off-by-one] A bug where a number is one too big or one too small, like using \texttt{len} as the last index. {\small\color{kcgray}(week 3)}
\item[open / with] \texttt{open} gets a file ready to use; \texttt{with} closes it again automatically. {\small\color{kcgray}(week 7)}
\item[pitch] a short talk that explains what your project is and why it is interesting. {\small\color{kcgray}(week 12)}
\item[polish] small finishing touches that make a working program clearer and nicer to use. {\small\color{kcgray}(week 11)}
\item[read / write / append] The modes \texttt{"r"}, \texttt{"w"} and \texttt{"a"}: look, replace everything, or add to the end. {\small\color{kcgray}(week 7)}
\item[showcase] an event where makers show their work and explain it. {\small\color{kcgray}(week 12)}
\item[version] one stage of a project; version 2 is the next, improved one. {\small\color{kcgray}(week 12)}
\end{description}
\end{multicols}

\end{document}
